\documentclass[11pt,a4paper]{article}
\usepackage[margin=1in]{geometry}
\usepackage{amsmath,amssymb,amsthm}
\usepackage{algorithm}
\usepackage{algpseudocode}
\usepackage{booktabs}
\usepackage{hyperref}

\theoremstyle{definition}
\newtheorem{definition}{Definition}[section]
\newtheorem{theorem}{Theorem}[section]
\newtheorem{lemma}[theorem]{Lemma}
\newtheorem{remark}{Remark}[section]

\title{\textbf{Rigorous Mathematical Specification, Geometric Margin Dynamics, and Convergence Analysis of the Perceptron Learning Rule}}
\author{Team Scaibu \\ \small Email: \texttt{jaydeep.9664920749@gmail.com} \\ \small Architecture and Core Engine Team}
\date{October 2026}

\begin{document}
\maketitle

\begin{abstract}
This document provides the foundational mathematical specification, geometric margin dynamics, and convergence proof for Rosenblatt's Perceptron Learning Algorithm. Operating on linearly separable binary data, the perceptron updates an affine hyperplane iteratively upon misclassification. We establish Novikoff's convergence theorem with explicit step-by-step proofs, detail the algorithmic pseudo-code matching the reference implementation, derive computational complexity, calculate a complete numerical example, and examine numerical edge cases.
\end{abstract}

\section{Notation and Mathematical Preliminaries}

Let $\mathcal{D} = \{(x_i, y_i)\}_{i=1}^N$ denote a training dataset of $N$ samples, where $x_i \in \mathbb{R}^d$ is the feature vector and $y_i \in \{0, 1\}$ is the binary target class label.

\begin{itemize}
    \item $y_i' \triangleq 2y_i - 1 \in \{-1, +1\}$: Bipolar signed label representation.
    \item $w \in \mathbb{R}^d$: Trainable weight vector of dimension $d$.
    \item $b \in \mathbb{R}$: Trainable scalar bias.
    \item $\tilde{x}_i \triangleq [x_i^T, 1]^T \in \mathbb{R}^{d+1}$: Homogeneous augmented input vector.
    \item $\tilde{w} \triangleq [w^T, b]^T \in \mathbb{R}^{d+1}$: Homogeneous augmented weight vector.
    \item $\eta > 0$: Learning rate step size.
    \item $R \triangleq \max_{i=1, \dots, N} \|x_i\|_2$: Maximum sample radius in feature space.
    \item $\gamma > 0$: Geometric separation margin.
\end{itemize}

\begin{table}[h]
\centering
\caption{Summary of Mathematical Symbols for Perceptron Learning}
\begin{tabular}{lll}
\toprule
\textbf{Symbol} & \textbf{Type / Domain} & \textbf{Description} \\
\midrule
$x_i$ & $\mathbb{R}^d$ & Input feature vector \\
$y_i$ & $\{0, 1\}$ & Binary label \\
$y_i'$ & $\{-1, +1\}$ & Signed bipolar target label ($2y_i - 1$) \\
$w, b$ & $\mathbb{R}^d, \mathbb{R}$ & Weight vector and bias term \\
$\eta$ & $\mathbb{R}_{> 0}$ & Learning rate \\
$R$ & $\mathbb{R}_{> 0}$ & Maximum sample Euclidean radius \\
$\gamma$ & $\mathbb{R}_{> 0}$ & Geometric margin \\
\bottomrule
\end{tabular}
\end{table}

\section{Problem Definition and Core Concepts}

\subsection{Intuition and Motivation}
The Perceptron is the foundational linear classification unit of neural networks. It defines an oriented decision hyperplane $\mathcal{H} = \{x \in \mathbb{R}^d : w^T x + b = 0\}$ dividing feature space into two half-spaces. If an incoming sample is correctly classified ($y_i'(w^T x_i + b) > 0$), no update occurs. If misclassified ($y_i'(w^T x_i + b) \le 0$), the hyperplane normal is rotated towards the misclassified point by adding $\eta y_i' x_i$, strictly reducing future error.

\subsection{Formal Mathematical Definition}
\begin{definition}[Linear Separability and Perceptron Update]
A dataset $\mathcal{D}$ is linearly separable with margin $\gamma > 0$ if there exists an optimal unit vector $u^* \in \mathbb{R}^{d+1}$ ($\|u^*\|_2 = 1$) such that:
\begin{equation}
y_i' (u^{*T} \tilde{x}_i) \ge \gamma \quad \forall i \in \{1, \dots, N\}
\end{equation}
For any sample $i$ where $\hat{y}_i' = \operatorname{sign}(w^T x_i + b) \ne y_i'$, the update rule is:
\begin{equation}
w \gets w + \eta y_i' x_i, \quad b \gets b + \eta y_i' \label{eq:perceptron_update}
\end{equation}
\end{definition}

\section{Algorithmic Specification}

The operational execution implemented in \texttt{NnAlgoPerceptronLearning} is shown below.

\begin{algorithm}[H]
\caption{Perceptron Learning Algorithm}
\begin{algorithmic}[1]
\Require Features $X \in \mathbb{R}^{N \times d}$, labels $Y \in \{0, 1\}^N$, learning rate $\eta > 0$, max epochs $E \ge 1$
\Ensure Learned weights $w \in \mathbb{R}^d$, bias $b \in \mathbb{R}$, convergence flag, epoch and update counts
\State $w \gets \mathbf{0} \in \mathbb{R}^d, \quad b \gets 0.0$
\State $\text{converged} \gets \textbf{false}, \quad \text{total\_updates} \gets 0$
\For{$\text{epoch} \gets 1 \textbf{ to } E$}
    \State $\text{errors} \gets 0$
    \For{$i \gets 1 \textbf{ to } N$}
        \State $y_i' \gets 2 \cdot Y[i] - 1$
        \State $\text{score} \gets \sum_{j=1}^d w[j] \cdot X[i][j] + b$
        \State $\hat{y} \gets \begin{cases} 1 & \text{if } \text{score} > 0 \\ 0 & \text{otherwise} \end{cases}$
        \If{$\hat{y} \ne Y[i]$}
            \For{$j \gets 1 \textbf{ to } d$}
                \State $w[j] \gets w[j] + \eta \cdot y_i' \cdot X[i][j]$
            \EndFor
            \State $b \gets b + \eta \cdot y_i'$
            \State $\text{errors} \gets \text{errors} + 1$
            \State $\text{total\_updates} \gets \text{total\_updates} + 1$
        \EndIf
    \EndFor
    \If{$\text{errors} = 0$}
        \State $\text{converged} \gets \textbf{true}$
        \State \textbf{break}
    \EndIf
\EndFor
\State \Return $\{w, b, \text{converged}, \text{epochs\_trained}: \text{epoch}, \text{total\_updates}, \text{final\_error\_count}: \text{errors}\}$
\end{algorithmic}
\end{algorithm}

\section{Theoretical Guarantees and Novikoff's Convergence Theorem}

\begin{theorem}[Novikoff's Perceptron Convergence Theorem (1962)]
Let $\mathcal{D} = \{(x_i, y_i')\}_{i=1}^N$ be linearly separable with augmented margin $\gamma > 0$ under unit vector $u^*$. Let $R = \max_i \|\tilde{x}_i\|_2$. Starting from $\tilde{w}_0 = \mathbf{0}$, the perceptron makes at most $k_{\max}$ mistake updates:
\begin{equation}
k_{\max} \le \left(\frac{R}{\gamma}\right)^2
\end{equation}
\end{theorem}

\begin{proof}
Let $\tilde{w}_k$ denote the augmented weight vector after $k$ updates.
\textbf{Lower Bound}: Using the update $\tilde{w}_k = \tilde{w}_{k-1} + \eta y_i' \tilde{x}_i$:
\begin{equation}
u^{*T} \tilde{w}_k = u^{*T} \tilde{w}_{k-1} + \eta y_i' (u^{*T} \tilde{x}_i) \ge u^{*T} \tilde{w}_{k-1} + \eta \gamma
\end{equation}
By induction from $\tilde{w}_0 = \mathbf{0}$, $u^{*T} \tilde{w}_k \ge k \eta \gamma$. By Cauchy-Schwarz, $\|\tilde{w}_k\|_2 = \|u^*\|_2 \|\tilde{w}_k\|_2 \ge u^{*T} \tilde{w}_k \ge k \eta \gamma$, so $\|\tilde{w}_k\|_2^2 \ge k^2 \eta^2 \gamma^2$.

\textbf{Upper Bound}: Since an update occurs only when $y_i' (\tilde{w}_{k-1}^T \tilde{x}_i) \le 0$:
\begin{equation}
\|\tilde{w}_k\|_2^2 = \|\tilde{w}_{k-1}\|_2^2 + 2 \eta y_i' (\tilde{w}_{k-1}^T \tilde{x}_i) + \eta^2 \|\tilde{x}_i\|_2^2 \le \|\tilde{w}_{k-1}\|_2^2 + \eta^2 R^2
\end{equation}
By induction, $\|\tilde{w}_k\|_2^2 \le k \eta^2 R^2$. Combining both bounds:
\begin{equation}
k^2 \eta^2 \gamma^2 \le \|\tilde{w}_k\|_2^2 \le k \eta^2 R^2 \implies k \le \left(\frac{R}{\gamma}\right)^2
\end{equation}
\end{proof}

\subsection{Limitations}
\begin{itemize}
    \item \textbf{Linear Separability Requirement}: If data is not linearly separable (e.g., XOR problem), the perceptron algorithm never converges and cycles indefinitely.
\end{itemize}

\section{Complexity Analysis}

\subsection{Time Complexity}
\begin{itemize}
    \item \textbf{Per Epoch}: Evaluating $N$ samples of dimension $d$ requires $\mathcal{O}(Nd)$ FLOPs.
    \item \textbf{Total Time}: $\mathcal{O}(E \cdot N d)$ where $E \le \lceil R^2 / \gamma^2 \rceil$.
\end{itemize}

\subsection{Space Complexity}
\begin{itemize}
    \item \textbf{Auxiliary Memory}: Stores $w \in \mathbb{R}^d$ and $b \in \mathbb{R}$, requiring $\Theta(d)$ auxiliary space.
\end{itemize}

\section{Worked Numerical Example}

Consider $N=2, d=2$:
$x_1 = [1.0, 2.0]^T, y_1 = 1 \implies y_1' = +1$.
$x_2 = [-1.0, -1.0]^T, y_2 = 0 \implies y_2' = -1$.
Let $\eta = 1.0, w_0 = [0.0, 0.0]^T, b_0 = 0.0$.

\begin{enumerate}
    \item \textbf{Sample 1}: $\text{score} = 0(1) + 0(2) + 0 = 0 \implies \hat{y}_1 = 0 \ne 1$. Update:
    \begin{align}
    w &\gets [0, 0]^T + 1.0(+1)[1, 2]^T = [1.0, 2.0]^T \\
    b &\gets 0 + 1.0(+1) = 1.0
    \end{align}
    \item \textbf{Sample 2}: $\text{score} = 1(-1) + 2(-1) + 1 = -2 \implies \hat{y}_2 = 0 == y_2$. Correct!
    \item \textbf{Epoch 2 Verification}:
    Sample 1: $\text{score} = 1(1) + 2(2) + 1 = 6 > 0 \implies \hat{y}_1 = 1$.
    Sample 2: $\text{score} = -2 < 0 \implies \hat{y}_2 = 0$.
    Zero errors $\implies$ Converged in 2 epochs with 1 update.
\end{enumerate}

\section{Numerical Considerations and Edge Cases}

\begin{itemize}
    \item \textbf{Zero-Crossing Decision Threshold}: Strict inequality ($w^T x + b > 0$) is used to ensure deterministic binary output assignment.
    \item \textbf{Scale Invariance}: Scaling all inputs by scalar $c > 0$ scales $R \to c R$ and $\gamma \to c \gamma$, keeping $k_{\max} \le (R/\gamma)^2$ strictly invariant.
\end{itemize}

\begin{thebibliography}{9}
\bibitem{rosenblatt1958} F.~Rosenblatt, ``The perceptron: A probabilistic model for information storage and organization in the brain,'' \emph{Psychological Review}, vol.~65, no.~6, pp.~386--408, 1958.
\bibitem{novikoff1962} A.~B.~Novikoff, ``On convergence proofs on perceptrons,'' in \emph{Symposium on the Mathematical Theory of Automata}, vol.~12, pp.~615--622, 1962.
\end{thebibliography}

\end{document}
